\section{The parameter complete intersections}
\label{sec:parameters}

The constraints from Lemma~\ref{lem:simple-factor} will be imposed on a
sum of many powers.  Their parameter space must remain integral, and its
smooth cone must have enough connectivity to support a later equivariant
construction.  We prove both properties from the multiplicity-one factor
condition.  All homology and cohomology in this section have integral
coefficients.

The join statement below is the global homogeneous form of the
Thom--Sebastiani phenomenon; compare
\cite[Proposition~4.1 and Appendix]{OhmotoYokura} for a local
nonisolated version. We give the needed global homotopy argument
directly, retaining integral homology throughout.

\begin{lemma}[Nonzero levels and joins]
\label{lem:homogeneous-joins}
Let \(w\) be a nonzero homogeneous polynomial of degree \(b\ge2\) on a
complex vector space \(U\).  Suppose that some irreducible factor of \(w\)
has multiplicity one.  Then \(w^{-1}(1)\) is connected.  For every
\(M\ge1\), the smooth hypersurface
\[
 F_w=\left\{(\ell_1,\ldots,\ell_M)\in U^{\oplus M}:
                 \sum_{m=1}^M w(\ell_m)=1\right\}
\]
has the homotopy type of the join of \(M\) copies of \(w^{-1}(1)\).
Consequently,
\[
 \widetilde H_j(F_w;\Z)=0\qquad(0\le j\le2M-2).
\]
\end{lemma}

\begin{proof}
Projectivization makes \(w^{-1}(1)\) a cyclic cover of degree \(b\) of
\(\PP(U)\setminus\{w=0\}\).  This complement is connected.  A small
meridian around a smooth point of the multiplicity-one factor, away
from all other factors, permutes the \(b\) choices of a root by a
generator of \(\mu_b\).  The monodromy is therefore transitive, proving
connectedness.

For the join assertion, consider nonzero homogeneous polynomials
\(f\) and \(g\), of positive degrees \(a\) and \(c\), in disjoint sets
of variables.  Cover the level \(f(x)+g(y)=1\) by
\[
 V_-=\{\operatorname{Re}f(x)<2/3\},
 \qquad
 V_+=\{\operatorname{Re}f(x)>1/3\}.
\]
On \(V_-\), a continuous \(c\)-th root of \(1-f(x)\) identifies
\(V_-\) with
\[
 f^{-1}\bigl(\{\operatorname{Re}z<2/3\}\bigr)\times g^{-1}(1).
\]
The first factor contracts to the origin by \(x\mapsto tx\),
\(0\le t\le1\), because \(f(tx)=t^a f(x)\).  Similarly,
\(V_+\simeq f^{-1}(1)\).  On the intersection both \(f(x)\) and
\(1-f(x)\) admit continuous roots, giving a homeomorphism
\[
 V_-\cap V_+\simeq
 f^{-1}(1)\times g^{-1}(1)\times
 \{z\in\C:1/3<\operatorname{Re}z<2/3\}.
\]
Under these identifications, the two inclusion maps are homotopic
to the projections onto \(g^{-1}(1)\) and \(f^{-1}(1)\).
The homotopy pushout for this open cover is their join.  Iteration
proves the asserted description of \(F_w\).  The integral homology
formula for joins, including its Tor terms, now gives the stated
vanishing: joining \(M\) connected spaces has no reduced homology
below degree \(2M-1\).  Smoothness of every nonzero homogeneous level
follows from Euler's identity.
\end{proof}

\begin{theorem}[The parameter spaces]
\label{thm:parameters}
Let \(k\) be a field admitting an embedding into \(\C\), let \(U\) be a
finite-dimensional \(k\)-vector space with \(\dim U\ge3\), and let
\[
 W\subseteq\Sym^b(U^\vee),\qquad b\ge2,\qquad s=\dim W.
\]
Assume that every nonzero element of \(W_{\bar k}\) has an irreducible
factor of multiplicity one over an algebraic closure \(\bar k\).
For integers \(h\ge2\) and \(M\ge2h+4s+10\), put \(N_*=M\dim U\)
and define
\[
 \begin{split}
 Y&=\left\{(\ell_1,\ldots,\ell_M)\in U^{\oplus M}:
                 \sum_{m=1}^M w(\ell_m)=0\quad(w\in W)\right\},\\
 X&=\PP(Y)\subseteq\PP(U^{\oplus M}).
 \end{split}
\]
Then \(X\) is a geometrically integral complete intersection of
codimension \(s\), cut out by \(s\) equations of degree \(b\), and
\[
 \operatorname{codim}_X\operatorname{Sing}X
       \ge M-2s+1>h,
\]
with the codimension of the empty locus understood to be infinite.
In particular, a general \(h\)-dimensional projective linear section
is a smooth complete intersection contained in \(X_{\sm}\).

After any embedding \(k\hookrightarrow\C\), let \(Y^o\) be the
nonzero cone over \(X_{\sm}\).  Both \(Y^o\) and its unit-norm
retract are \((M+2h)\)-connected.
\end{theorem}

\begin{proof}
All geometric assertions may be checked after extension to \(\C\).
The multiplicity-one hypothesis persists under this extension.
Indeed forms without a multiplicity-one factor constitute a finite
union of images of projective multiplication maps
\((p_1,\ldots,p_t)\mapsto\prod_\nu p_\nu^{a_\nu}\), with
\(a_\nu\ge2\) and \(\sum_\nu a_\nu\deg p_\nu=b\).
Their intersection with \(\PP(W)\) is empty geometrically over \(k\),
and remains empty after extension.
Write \(r=\dim U\), \(d=N_*-s\), and \(\kappa=M+2h\).
If \(s=0\), then \(X=\PP^{N_*-1}\) and
\(Y^o=\C^{N_*}\setminus\{0\}\simeq S^{2N_*-1}\).
The conclusions follow from \(2N_*-2\ge\kappa\).
Henceforth assume \(s\ge1\).

\emph{Dimension and the singular locus.}
Choose a basis \(w_1,\ldots,w_s\) of \(W\).  A point has defective
Jacobian rank for the defining equations of \(Y\) only if, for some
\([w]\in\PP(W)\),
\[
 \nabla w(\ell_1)=\cdots=\nabla w(\ell_M)=0.
\]
For any nonzero \(w\), its gradient vanishing set has dimension at
most \(r-1\).  The incidence with \(\PP(W)\) therefore bounds the
defective-rank locus in \(U^{\oplus M}\) by
\begin{equation}
 \dim\{\text{defective rank}\}\le N_*-M+s-1.
 \label{eq:parameter-jacobian}
\end{equation}
Every irreducible component of \(Y\) has dimension at least \(N_*-s\).
Since \(M>2s-1\), no component is contained in the locus in
\eqref{eq:parameter-jacobian}.  Full Jacobian rank gives local dimension
\(N_*-s\), so all components have dimension \(d\) and are generically
smooth.  The \(s\) equations thus have height \(s\) in the polynomial
ring and form a regular sequence.  The resulting complete intersection
is Cohen--Macaulay, with no embedded associated primes.  As it is
generically reduced, it is reduced.

Away from the origin, the cone projection is locally a product with
\(\C^\times\), so smoothness of \(Y\) is equivalent to smoothness of
\(X\).  Equation~\eqref{eq:parameter-jacobian} gives
\begin{equation}
 \operatorname{codim}_Y(Y\setminus Y^o)\ge c:=M-2s+1.
 \label{eq:parameter-omitted}
\end{equation}
Here the origin causes no additional restriction: \(d\ge c\).
The same bound applies to \(\operatorname{Sing}X\).
It remains to prove integrality and connectivity; we have not yet
assumed that the smooth locus is connected.

\emph{Homology of the complement.}
A point outside \(Y\) makes some linear combination of the defining
equations nonzero.  Recording that combination replaces the complement
by a space with contractible fibers over it, while exposing the
nonzero hypersurface levels of Lemma~\ref{lem:homogeneous-joins}.
Define the open incidence variety
\[
 \cH=\left\{(\ell,[w])\in\C^{N_*}\times\PP(W):
                         \sum_m w(\ell_m)\ne0\right\}.
\]
Projection to \(\C^{N_*}\setminus Y\) is an affine bundle of rank
\(s-1\).  For a fixed \(\ell\notin Y\), each line \([w]\) in the
fiber has a unique representative satisfying \(\sum_m w(\ell_m)=1\).
Thus the fiber is the affine hyperplane on which the nonzero evaluation
functional on \(W\) equals one.
Consequently this projection is a homotopy equivalence.

Evaluation also defines a map
\[
 p:\cH\longrightarrow
 \mathcal B:=\cO_{\PP(W)}(1)^\times.
\]
The value at \((\ell,[w])\) is the nonzero linear functional on the
tautological line \(kw\) obtained by evaluation at \(\ell\).
Euler's identity shows that \(p\) is a smooth submersion.
Its fibers have complex dimension \(N_*-1\); after choosing a
representative of \([w]\) and rescaling, they are the levels
\(F_w\) of Lemma~\ref{lem:homogeneous-joins}.

The join lemma controls the homology of every fiber of \(p\).
To pass from these individual fibers to \(\cH\), we use direct image
with proper support; we do not require \(p\) to be a locally trivial
fibration.  Fiber homology becomes compact-support cohomology near
the top degree, where the Leray spectral sequence gives the needed
comparison.  Set \(q=2N_*-2\).  Compact-support Poincar\'e duality on
the fibers \cite[Theorem~3.35]{HatcherAT} and
Lemma~\ref{lem:homogeneous-joins} give
\begin{equation}
 R^q p_!\Z=\Z,\qquad
 R^i p_!\Z=0\quad
 \begin{cases}
 i>q,\\
 q-(2M-2)\le i<q.
 \end{cases}
 \label{eq:parameter-support-rows}
\end{equation}
The top sheaf is constant: a submersion
chart over a small base open set supplies an oriented open ball
in each fiber.  Extension by zero maps its top compact-support
orientation class to the positive generator for that fiber.
This gives an isomorphism from the constant sheaf to the top
direct-image sheaf on the open set.  The local isomorphisms agree
by the complex orientation.

The compact-support Leray spectral sequence is
\[
 H_c^a(\mathcal B,R^i p_!\Z)
       \ \Longrightarrow\ H_c^{a+i}(\cH;\Z).
\]
The stalk formula and this spectral sequence follow from proper-support
base change and derived composition
\cite[Theorems~5.9 and~5.10]{SchnurerSoergel16}; maps between the locally
compact Hausdorff spaces here satisfy their hypotheses. Equivalently,
one may compactify \(p\), extend the constant sheaf by zero, and apply
proper base change and Leray. The base is a real \(2s\)-dimensional
manifold, so its compact-support cohomology with coefficients in any
abelian sheaf vanishes above degree \(2s\)
\cite[Proposition~5.1.2(ii)]{Schapira}. In total degree \(q+2s-j\), with
\(0\le j\le2M-2\), only the top row of
\eqref{eq:parameter-support-rows} can contribute.  In fact, the next
possible row is \(i=q-(2M-1)\), which would require
\[
 a\ge2s-j+2M-1>2s.
\]
An outgoing differential from the top row to another nonzero row
has length at least \(2M>2s\), and hence has zero target.  An
incoming differential would come from above the top row.
Poincar\'e duality on \(\cH\) and \(\mathcal B\) therefore yields
\begin{equation}
 H_j(\cH;\Z)\cong H_j(\mathcal B;\Z)
       \qquad(0\le j\le2M-2).
 \label{eq:parameter-complement-homology}
\end{equation}
The punctured line bundle \(\mathcal B\) retracts to
\(S^{2s-1}\).

\emph{Homology of the smooth cone.}
The closed/open compact-support sequence for
\(Y\subseteq\C^{N_*}\), followed by duality on its complement,
gives
\begin{equation}
 H_c^{2d-j}(Y;\Z)
   \cong H_{2s+j-1}(\C^{N_*}\setminus Y;\Z)
   \cong H_{2s+j-1}(\cH;\Z)
 \label{eq:parameter-support-transfer}
\end{equation}
for \(0\le j\le\kappa\).  Both adjacent ambient compact-support
groups vanish, since their degrees are below \(2N_*\).
The homology range in \eqref{eq:parameter-complement-homology}
applies because
\begin{equation}
 2s+\kappa-1\le2M-2.
 \label{eq:parameter-complement-range}
\end{equation}
Removing \(Y\setminus Y^o\) does not change the left side of
\eqref{eq:parameter-support-transfer}. Indeed, a finite filtration of
the removed algebraic set by smooth locally closed strata, together with
the compact-support closed/open sequence, gives vanishing above its real
dimension from the same manifold bound.
Equation~\eqref{eq:parameter-omitted} bounds that dimension by
\(2(d-c)\), whereas
\begin{equation}
 2d-j-1>2(d-c)\quad(0\le j\le\kappa),
 \qquad
 \kappa\le2c-2=2M-4s.
 \label{eq:parameter-removal-range}
\end{equation}
Both \eqref{eq:parameter-complement-range} and
\eqref{eq:parameter-removal-range} follow from
\(M\ge2h+4s+10\).
Duality on \(Y^o\) now gives
\[
 H_0(Y^o;\Z)=\Z,\qquad
 \widetilde H_j(Y^o;\Z)=0\quad(1\le j\le\kappa).
\]
Every irreducible component of \(Y\) meets \(Y^o\) densely.
The smooth loci of distinct components are disjoint, so their
union could be connected only if there is one component.
Thus \(Y\), and hence \(X\), is integral.

\emph{From homology to homotopy.}
The preceding argument leaves only simple connectivity to prove.
The quasi-projective Lefschetz theorem
\cite[Theorem~1.1.3(ii)]{HammLe85} applies to the smooth open
\(X_{\sm}\) in the projective variety \(X\), with deleted algebraic
subset \(\operatorname{Sing}X\).
It states that the smooth open has the homotopy type of a space
obtained from a sufficiently general hyperplane section by
attaching cells of dimension at least the complex dimension
of the open.  Apply the theorem successively, intersecting the deleted
subset with each hyperplane as well.  Each section down to a surface
therefore surjects on the fundamental group of the preceding smooth
open.

Since \(\operatorname{codim}_X\operatorname{Sing}X>h\ge2\),
a general surface linear section avoids the singular locus.
It is a smooth complete intersection \(S_0\subseteq\PP^{s+2}\)
of \(s\) hypersurfaces of degree \(b\).  Such a surface is simply
connected.  To check the usual smooth flag needed here, replace
its defining equations by general invertible linear combinations.
Off their common base locus, Bertini gives smooth successive
intersections.  Along \(S_0\), the independent differentials of
the defining equations give smoothness of every initial
intersection.  Smooth projective Lefschetz, applied to these
very ample divisors, preserves the fundamental group down to
dimension two, starting from projective space.  Therefore
\(\pi_1(S_0)=0\), and the surjections above give
\(\pi_1(X_{\sm})=0\).

The cone \(Y^o\) is the complement of the zero section in
\(\cO_{X_{\sm}}(-1)\).  Its unit circle bundle has a homotopy
exact sequence in which
\(\pi_1(S^1)=\Z\) surjects onto \(\pi_1(Y^o)\), since the base
is simply connected.  Hence \(\pi_1(Y^o)\) is abelian.
Its abelianization is \(H_1(Y^o;\Z)=0\), so it is trivial.
Smooth complex algebraic varieties have CW type.  Successive
applications of the Hurewicz theorem \cite[Theorem~4.32]{HatcherAT}
to the homology vanishing
already proved now give \(\pi_j(Y^o)=0\) for \(1\le j\le\kappa\).

Radial retraction gives the same connectivity for the unit cone.
Finally, the strict bound on the singular codimension and Bertini
give the asserted smooth \(h\)-dimensional sections.
All conclusions hold after every complex embedding and imply the
stated geometric assertions over \(k\).
\end{proof}

We shall use the unit cone
\[
 Y_1=\{y\in Y^o:\|y\|=1\},
\]
with a Hermitian norm that is the orthogonal sum of norms on the
\(M\) blocks.  Common circle scaling and independent multiplication
of the blocks by \(b\)-th roots of unity preserve \(Y\) and its
smooth locus.  The radial retraction \(Y^o\to Y_1\) is equivariant
for these commuting actions.
